\section{Experimental design}
\label{sec:experiments}

\subsection{Offline model-selection experiment}

The selection experiment is $5$ variants $\times$ $3$ seeds, giving 15
validation runs. Primary comparison: rollout-macro NLL. Secondary comparisons:
top-1 ADE/FDE, reactive-subset ADE and latency. A one-shot frozen test evaluates
the validation-selected checkpoint of each variant and the pretrained \Bzero{}
control. This tests H1 and H3 without test-set tuning.

\subsection{Sequence-use and calibration diagnostics}

For \Tone{} and \Ttwo{}, valid interaction tokens are replaced by zeros or
shuffled while weights remain frozen. An increase in NLL supports input use,
but not causal interaction understanding. Aggregate and response-active
calibration are reported separately. These analyses test H2 and diagnose
tail behaviour; they cannot reselect the model.

\subsection{Deployment reliability gate}

Day9 runs eight smoke arms through raster/sequence input, GMM decoding,
calibration, risk allocation, solver and supervisor. Passing this gate supports
deployability but is not an effect experiment. Formal closed-loop conclusions
are based only on Days10--12.

\subsection{Nominal closed-loop factorial}

The nominal design is
\[
2\ \text{predictors}\times4\ \text{risk policies}\times
2\ \text{target styles}\times5\ \text{initialisations}=80
\]
rollouts. Predictors are \Bzero{} and frozen \Bone{}. Risk policies are fixed
aggressive, fixed medium, fixed conservative and adaptive. The matrix evaluates
the full fixed frontier rather than a single hand-picked control.

\subsection{Arrival-timing robustness and pooled synthesis}

The shifted-timing design is
\[
2\ \text{predictors}\times2\ \text{risk policies}\times
2\ \text{target styles}\times2\ \text{offsets}\times
5\ \text{initialisations}=80
\]
rollouts, using fixed medium/adaptive risk and target offsets $-3$ and $+3$~m.
Combining these with matching nominal ($0$~m) conditions yields 24 cells and
120 formal rollouts for the three-offset synthesis. This design tests H4--H6
and records delay, footprint separation, solver failure, supervisor activity
and observed reliability events.

\subsection{Controls and validity safeguards}

Key controls are the pretrained \Bzero{}, simple \Bone{}, matched MLP adapters,
fixed-risk frontier, target-style factorial, paired initialisations, train-only
normalisation, validation-only selection, frozen calibration/deployment hashes,
and a collision-filtered sensitivity. Table~\ref{tab:threats} later separates
each mitigation from the limitation that remains.

